\section{Introduction}
\label{sec:intro}
%
Recent advancements in large language models (LLMs) such as GPT-5~\cite{gpt5,openai2023gpt4}, DeepSeek-R1~\cite{guo2025deepseek}, LLaMA-4~\cite{touvron2023llama, llama3,llama4} and QWen 3~\cite{qwen2025qwen25technicalreport,yang2025qwen3technicalreport}, have shown strong capabilities of generating natural language texts across various application domains~\cite{zhang2019pretraining, liu2021makes, stiennon2020learning,dosovitskiy2020image}.
%
Given the increasingly high cost of LLM pre-training, there is a trend of standardizing foundational pre-trained LLMs and sharing sub-models across multiple downstream tasks~\cite{houlsby2019parameter,zhang2021share} through finetuning (i.e., continuous training on small task-specific datasets). For example, a series of {\em parameter-efficient finetuning} (\peft) techniques~\cite{houlsby2019parameter,li2021prefix,pfeiffer2020adapterfusion,hu2021lora,he2021towards,lester2021power} have been proposed to update a subset of trainable parameters from an LLM or introduce a small number of new parameters into the LLM while keeping the vast majority of the original LLM parameters frozen. The base LLM sharing paradigm provides opportunities to multiplex resources and improve cluster utilization (e.g., Google~\cite{barham2022pathways}, Microsoft~\cite{hu2021lora}, and Tencent~\cite{nie2023angel}).
Many recent efforts~\cite{zhou2022pets,wu2024dlora,sheng2023s,chen2023punica,zhao2024lora,iliakopoulou2024chameleon} inspired by these opportunities to build multi-task serving systems, handling inference requests for multiple fine-tuned LLMs simultaneously in a single GPU cluster.

In addition to inference, almost all LLM serving companies (including OpenAI GPT-4o~\cite{openaifinetuning}, Google Gemini~\cite{googlefinetuning}, and Databricks Mosaic AI~\cite{databricksfinetuning}), now offer commercial finetuning APIs to facilitate the creation of custom models.
This trend is not only due to the widespread adoption of LLMs but also driven by the growing demand for task-specific model customization~\cite{liu2024customizing}, data access control~\cite{tan2024democratizing}, and the need for continuous and frequent model updating~\cite{gao2023unified}.

Currently, the most popular way to offer both inference and finetuning services is to run each task on separate, {\em dedicated} clusters~\cite{weng2022mlaas,gao2022deep,chen2023deepboot}, which exclusively use hardware resources through specialized systems optimized primarily for latency and throughput, respectively. While straightforward, this practice leads to significant economic inefficiencies. In particular, resources must be overprovisioned for inference tasks, as real-world LLM inference workloads exhibit highly unpredictable, bursty request arrival patterns~\cite{wang2024towards}. For example, public reports from Microsoft and Alibaba cite average GPU utilization rates of just 52\%~\cite{lithos2025} and 10\%~\cite{lithos2025}, respectively, with production services showing similar underutilization patterns. While inference tasks must be provisioned for peak burst capacity to maintain SLOs, the resulting idle resources during normal operation cannot be utilized by finetuning tasks, which could tolerate the latency variations (minutes to hours vs. milliseconds). This is economically unsustainable given the high cost and power demands of modern GPUs, which can now exceed 1,000W per chip~\cite{lithos2025}.

In this paper, we explore a key research question: {\em Can we design a system that {\em simultaneously} serves both inference and finetuning tasks within a shared GPU cluster, dynamically adapting to unpredictable burst patterns while maintaining strict inference SLOs?}
Previous attempts (\Cref{sec:coserving}) at multiplexing inference and finetuning tasks, while successful in increasing the GPU utilization, have seen limited adoption due to their impracticality.

To solve this fundamental challenge, we introduce \textit{co-serving}, a novel multiplexing technique that effectively handles bursty workloads while satisfying strict SLOs. We built \sys, the first system for LLM inference and finetuning that implements this technique. \sys introduces a \textit{PEFT-as-a-service (PaaS)} interface, which unifies inference and finetuning tasks, enabling their joint execution on shared GPU resources.
The key insight behind co-serving is that inference and finetuning tasks, when using the same base LLMs, can be merged at fine granularity—at the level of individual tokens rather than entire requests or kernels. This token-level scheduling enables millisecond-scale resource reallocation in response to inference bursts.

This co-serving approach offers critical advantages for handling bursty workloads. When inference requests suddenly spike, \sys can instantly throttle finetuning tokens within the same GPU kernel execution, reallocating resources to maintain inference SLOs without the overhead of context switching or reconfiguration. During normal load, finetuning tasks opportunistically consume the idle capacity reserved for bursts, improving overall utilization. Since the base LLM parameters are shared and frozen during finetuning, memory overhead is minimized, enabling efficient multiplexing.

Overall, this paper makes the following key contributions:
\begin{itemize}[leftmargin=10pt,parsep=1pt,topsep=2pt]
\item We introduce FlexLLM, the first system to co-serve LLM inference and PEFT finetuning on shared GPUs through token-level computation fusion.
\item We develop static optimizations that reduce GPU memory requirements by up to 80\% through dependent parallelization and graph pruning.
\item We propose token-level finetuning with hybrid scheduling that maintains strict latency SLOs while maximizing GPU utilization.
\item We achieve 1.9$-$4.8$\times$ finetuning throughput improvements under heavy loads and 2.5$-$6.8$\times$ under light loads.
\end{itemize}